\section{The Computational Reason for Bizarreness}\label{sec:obh_bizarreness}

A simple ``replay'' of data, however, creates a fundamental problem familiar to artificial intelligence: \textit{overfitting}. A neural network (or a brain) trained exclusively on the exact data of its past experiences becomes rigid. It may excel at those specific, known tasks but will fail to \textit{generalize} its knowledge to novel, unseen situations \citep{hoel2021overfitted}.

This has led to the ``overfitted brain hypothesis'', which proposes that the \textit{purpose} of dreaming \- particularly its bizarre nature \- is to \textit{prevent} this overfitting \citep{hoel2021overfitted}. According to this model, dreams function as ``augmented'' or ``corrupted'' versions of waking experience. By injecting this ``noise'' or ``counterfactuals'' into the system, the brain regularizes its own model, ensuring that its learned concepts remain robust, flexible, and generalizable.

This theory reframes the ``bizarreness'' of dreams as a critical functional feature, not a meaningless flaw. The brain is deliberately ``corrupting'' its own memories to stress-test its understanding of the world.
